
Self-supervised learning methods achieve up to 90\% worst-group accuracy on spurious correlation benchmarks when using linear probes on frozen embeddings. The geometric structure that enables this performance is not well characterized. A common assumption is that spurious features manifest as discrete, geometrically separable clusters measurable via clustering metrics such as Adjusted Mutual Information (AMI). We tested this assumption through controlled experiments on the Waterbirds dataset. SimCLR training (20 epochs, ResNet-50) produced embeddings with AMI=0.2795, below the 0.4 threshold conventionally used to indicate reliable clusterability. Learning-speed aware SSL (LA-SSL) yielded AMI=0.2852, representing a 2.04\% relative increase rather than the hypothesized 30\% reduction. Linear separability metrics remained high ($AUC\approx 0.98)$ despite low clusterability scores. These results indicate that linear separability and discrete clusterability are dissociable properties in SSL embeddings under the tested conditions. The experiments were conducted as proof-of-concept validation (20 epochs) prior to full-scale training; the mechanism hypotheses tested were not supported at this training duration. Implementation validation was performed for cluster-balanced retraining mechanisms (43/43 tests passing, 100\% software development compliance for h-e1; 5/5 tests passing, 100\% compliance for h-m-integrated), though experimental validation of intervention efficacy was not completed.
